\section{Discussion}
\label{sec:discussion}

\subsection{What the results establish}
\label{sec:what-established}

The central claim of the paper is an architectural separation: adaptation may
choose among safe actions, but it may not choose whether an action is safe. The
experiments support that separation in a specific and limited sense. The
invariant-violation rate is zero for every policy that enforces I1/I2, and it
remains zero in the scenario where the belief is demonstrably wrong. This is not
because the estimator is good; it is because the estimator has no authority over
the floor. The Risk-Only ablation makes the counterfactual concrete: the same
trust-weighted Bayesian machinery without the floor violates the floor at $327.41$
per $1{,}000$ decisions, close to the $393.31$ of a policy that ignores risk
entirely. The node-side PEP re-check is what closes the gap between a gateway
decision and an enforced action.

The utility result is positive but modest, and should be read carefully. Full
attains $0.120$ mean attack-scenario utility against $0.084$ for Static-High and
$0.077$ for Threshold, a paired gain of $+0.036$ and $+0.043$
($p_{\text{Holm}}=0.008$ in both cases). It also cuts the false-denial rate from
$0.088$ to $0.027$. Those are real effects on a normalised scale, but they are
smaller than the differences against policies that are either unsafe (Static-Low,
Risk-Only) or over-deny. A reader who expects a context-adaptive controller to
dominate a static-high posture on every metric will not find that here, and the
paper does not claim it: Static-High still has lower residual exposure
($0.127$ versus $0.177$).

The belief result is conditional, and the condition is now explicit. Trust-weighted
fusion is valuable when the adversary poisons a subset of modalities: the Brier
score under energy inflation improves from $0.455$ for uniform fusion to $0.003$
for Full. It is also sufficient for the in-model downgrade, where the attested
device anchor and the cross-node dispersion survive the spoof: the belief AUROC is
$1.000$ and the Brier score is $0.003$. It is not sufficient when the adversary
also controls the attestation anchor. In the \texttt{tee} scenario the Brier score
is $0.437$, the poisoning AUROC is $0.273$ (below chance) and the malicious-grant
rate rises to $0.473$. That negative result is the strongest argument for the
architecture, because a design that relied on detection would fail exactly where
this one remains safe.

\subsection{Limitations}
\label{sec:limitations}

\paragraph{Simulation, not deployment.} All performance numbers come from a
parameterised discrete-time simulator. The observation model, transition matrix
and attack processes are explicit and the streams are paired, but they are not a
real deployment. Latency, energy and footprint are analytical estimates from
datasheet-level operation costs, not measurements on the three hardware classes.
The relative comparison between postures is meaningful; the absolute numbers are
not predictions for a specific board. A hardware testbed is the obvious next
step.

\paragraph{A subverted anchor is not solved.} The paper does not claim that the
belief detects a compromise of the node's trusted execution environment. On the
contrary, \cref{tab:detection} shows that both the poisoning AUROC ($0.273$) and
the belief AUROC ($0.210$) are below chance in the \texttt{tee} scenario, and
\cref{tab:utility} shows that Full's attack-scenario utility drops to $0.058$
there---far below its usual level, even though it edges Static-High ($0.053$)
because both postures become cheap. The guarantee survives because I1/I2 are
unconditional, but the posture selected inside the envelope can still be too
cheap, which is why residual malicious grants rise to $0.473$. Closing that gap requires evidence the
adversary cannot fabricate---hardware attestation rooted outside the TEE, physical
invariants, or out-of-band confirmation---not a better classifier on the same
features. The in-model downgrade, by contrast, is handled: the belief stays
calibrated because the anchor and dispersion are retained.

\paragraph{The floor bounds the action, not the credential.} A malicious request
that holds a valid credential and arrives with a plausible context can be granted
an action that satisfies I1. The invariant layer ensures that such an action uses
strong authentication, AEAD, freshness and inspection; it does not ensure that the
request is benign. This is why Full's malicious-grant rate ($0.177$) is higher than
Static-High's ($0.127$). Most of the gap is the out-of-model \texttt{tee} scenario:
excluding it, Full is $0.128$ against $0.104$. The floor is a bound on the
\emph{mechanism}, not on the \emph{authorisation decision}. Reducing residual
exposure requires either a stronger unconditional floor
(\cref{sec:rq5-floor}) or an authorisation layer that is outside the scope of this
paper.

\paragraph{Trust anchors.} The device modality is treated as an attestation
anchor. If a node's trusted execution environment is subverted, the anchor is no
longer trustworthy and the trust weighting can be steered by the compromised node;
this is the \texttt{tee} scenario, deliberately reported rather than hidden. I1
still limits what that node can do to other physical processes, but it cannot
prevent the node from exercising credentials it legitimately holds. The same caveat
applies to a compromised gateway: the PDP is trusted in this model, and a gateway
compromise is out of scope.

\paragraph{Model and tuning.} The emission model, transition matrix and cost
weights are the defender's assumptions. In this simulator the likelihood is
evaluated with the same parameters that generate the observations, so the belief
is correctly specified by construction and the calibration and detection results
are an oracle-model upper bound; a defender with a misspecified model would do
worse. The thresholds in \cref{tab:params} were fixed before the final ten-seed
matrix was generated, but
they were selected using earlier runs on the same simulator; this leaves a risk of
optimistic bias that a pre-registered, held-out parameter search would remove. The
violation penalty $\rho=6.0$ is a reporting convention, not a controller
parameter; because the proposed controller never grants a below-floor action, its
utility is unaffected by $\rho$, while the unconstrained baselines are penalised
for violations they actually commit. The floor sensitivity (\cref{tab:floor}) and
the component ablation (\cref{tab:ablation}) are included partly to make the
dependence on those choices visible. The action lattice is finite ($243$ actions),
the deployment is $24$ nodes, and the Byzantine fraction is $10/24$; larger or
heterogeneous deployments may shift the quantitative trade-offs even if the
structural argument is unchanged.

\paragraph{Baselines.} The baselines are deliberately simple. Threshold is a
univariate point estimate on the same anomaly feature the adversary manipulates;
we raise it to the structural floor so that its IVR is not an artefact, but we do
not evaluate a multivariate or learned detector tuned to Full's false-positive
rate. Static-Med is a fixed mid-point and Static-High is the strongest uniform
posture. A learned detector could close part of the utility gap without the
belief, and a detector trained on the simulator's oracle-model features would look
better than one trained on real telemetry. The comparison therefore isolates the
architecture---belief, floor and PEP---rather than the best achievable detector.

\paragraph{Trace calibration is partial.} The IoT-23 capture is a honeypot
dataset; its absolute rates are not comparable to per-node application telemetry,
and we match only the benign-to-malicious mean-log-rate contrast ($1.689$, with
per-capture values $1.131$ and $2.570$). The simulator's network dispersion is
narrower than the trace ($0.50$ and $0.55$ versus $0.65$ and $0.86$), so no
distributional-shape claim is made. The capture does not exercise the controller,
the request process or the physical process, so it cannot validate the utility or
safety results. It constrains one prior, and we describe it as such.

\paragraph{Adversary strength.} The attacker is white-box but not an optimal
POMDP solver. The adaptive scenario increases the spoofed fraction over time and
then holds the anomaly just below the public point-estimate threshold, and the
resonant scenario pulses near the dwell period, but neither is a worst-case policy
computed against the defender's model. A stronger adversary could search for
belief trajectories that keep the controller at the cheapest admissible action for
longer; the floor would still hold, but the utility and exposure numbers would move
against the defender.

\subsection{Threats to validity}
\label{sec:validity}

\paragraph{Construct validity.} Mission utility is a weighted sum of value,
damage, denial cost, energy and latency, and it is reported after subtracting a
violation penalty $\rho$ for any granted action below the floor. The weights are
design choices, and a different operator could rank postures differently. We
mitigate this by reporting the components separately (utility, false denial,
malicious grant, safety denial) rather than only the composite, and by showing the
floor sensitivity. The IVR, false-denial and safety-denial metrics are direct
counts and do not depend on the utility weights or on $\rho$.

\paragraph{Internal validity.} All policies within a seed and scenario replay the
same observation and request stream, which removes stream-to-stream variance from
the comparison. The simulator is deterministic given the seed, and the unit tests
check the invariant mask, the belief normalisation, the trust bounds, the
fail-safe fallback and the paired-stream property. The main internal threat is
parameter tuning on the same scenario family, discussed above. A second threat is
that the stochastic detection model is shared by all policies, so a policy cannot
gain by exploiting a different mechanism model; this is intentional for
comparability but means the absolute detection rates are model artefacts.

\paragraph{External validity.} The results are for a 24-node deployment with three
hardware classes, one action lattice and eight scenarios. The qualitative
claim---a hard floor decouples safety from estimation error---is not specific to
these numbers, but the magnitude of the utility and exposure trade-offs is. The
IoT-23 calibration improves the realism of one prior without making the simulator a
trace replay.

\paragraph{Conclusion validity.} The paired tests use ten seeds and a conservative
non-parametric test, with each seed's seven attack scenarios aggregated into one
mean before testing. That aggregation is a deliberate summary, not seven
independent scenario tests: with ten pairs the smallest attainable two-sided
$p$-value is $0.00195$, so every $p_{\text{Holm}}=0.008$ in \cref{tab:stats} sits
at the resolution floor and should be read as sign consistency across seeds, not
as a precise significance level. We therefore lead with effect sizes and the
per-scenario breakdown, and we do not claim scenario-specific significance. The utility
difference against Static-High ($+0.036$) is small but reliably positive
($p_{\text{Holm}}=0.008$); the exposure difference ($+0.050$) is larger and goes
against the proposed policy. Both are reported.

\subsection{Design implications}
\label{sec:implications}

Three implications follow for practitioners. First, the safety floor should be
chosen as an explicit risk budget, not as a by-product of estimator tuning. The
floor sensitivity in \cref{tab:floor} shows how strengthening I1 trades exposure
against utility and energy; an operator can pick a point on that curve without
changing the controller. Second, the attestation anchor is load-bearing: if the
device modality is not rooted in hardware the adversary cannot rewrite, the
downgrade failure becomes easier to induce. Third, the PDP/PEP split matters: the
belief and the scorer can run on a gateway, but the admissible-set check and the
fail-safe action must run on the node, or a network adversary can bypass them.

\subsection{Ethics and reproducibility}
\label{sec:ethics}

The work uses a public network capture and a simulator; it involves no human
subjects and no private data. The attack scenarios are defensive stress tests, and
the paper does not release an exploit. The dual-use risk is the usual one for
access-control research: a description of how context poisoning defeats a CAAC
policy could be misused. We judge the defensive value to outweigh that risk, and
the design specifically aims to make the attack less effective by removing the
estimator's authority over the safety floor.

The simulator, the analysis scripts and the generated tables are deterministic
given the seeds and are organised so that every reported number can be regenerated
from the frozen matrix. The manuscript does not contain hand-typed experimental
values; the asset generator reads the CSV and writes both the tables and the prose
summary used for cross-checking.
